% ============================================================
%  Chapitre 1 — Introduction
% ============================================================
\chapter{Introduction}

\section{Contexte et motivation}

La \textbf{Représentation des Connaissances et Raisonnement} (\textsc{rcr}) est
une discipline fondamentale de l'intelligence artificielle dont l'objectif est de
formaliser les connaissances d'un domaine de manière à permettre à un système
informatique de raisonner automatiquement à partir de ces connaissances.

Contrairement à la programmation classique qui encode des procédures explicites,
la \textsc{rcr} permet de décrire \emph{ce que l'on sait} et de laisser un
moteur d'inférence déduire \emph{ce qui en découle}. Cette approche déclarative
est particulièrement adaptée aux domaines où les connaissances sont incomplètes,
incertaines ou sujettes à des exceptions.

Plusieurs formalismes ont été développés au fil des décennies, chacun avec ses
forces propres : certains sont monotones et garantissent la cohérence logique
(logique classique, logique de description, logique modale), d'autres sont
non-monotones et permettent de réviser les conclusions face à de nouvelles
informations (logique des défauts, réseaux sémantiques avec héritage à exceptions).

\section{Objectifs du projet}

Ce projet vise à :

\begin{enumerate}
  \item \textbf{Implémenter} cinq formalismes de représentation des connaissances
        à l'aide de bibliothèques Java établies.
  \item \textbf{Appliquer} chaque formalisme à deux domaines réels et concrets,
        permettant une comparaison directe de leur expressivité.
  \item \textbf{Démontrer} le phénomène de non-monotonie pour les formalismes
        qui le supportent.
  \item \textbf{Développer} une interface graphique unifiée permettant
        d'interagir avec chaque raisonneur de manière pédagogique.
\end{enumerate}

\section{Domaines d'application}

Deux domaines ont été choisis pour leur richesse sémantique et leur
complémentarité :

\begin{description}
  \item[Industrie musicale algérienne] — Artistes (Soolking, Feu, Tif, Flenn),
    labels (Universal, Believe), albums, certifications (or, platine),
    relations contractuelles, collaborations internationales.
    Ce domaine se prête naturellement à la modélisation de statuts
    (artiste signé, indépendant, libre de droits) et de hiérarchies
    (rappeur algérien $\sqsubseteq$ rappeur $\sqsubseteq$ artiste).

  \item[Système judiciaire] — Magistrats (juge, procureur), avocats,
    accusés (mineur, récidiviste, présumé innocent), infractions (crime, délit,
    contravention), tribunaux (cour d'assises, tribunal correctionnel,
    tribunal pour enfants).
    Ce domaine illustre parfaitement les exceptions et les conflits de règles
    (un mineur ayant commis un crime grave relève-t-il du tribunal des enfants
    ou de la cour d'assises ?).
\end{description}

\section{Structure du rapport}

\begin{description}
  \item[Chapitre 2] présente l'architecture technique du projet (organisation
        Maven, technologies, interface graphique).
  \item[Chapitres 3 à 7] détaillent chacun des cinq formalismes :
        fondements théoriques, implémentation, application aux deux domaines,
        résultats.
  \item[Chapitre 8] propose une comparaison transversale des formalismes.
  \item[Chapitre 9] conclut et ouvre des perspectives.
\end{description}
